\documentclass[lettersize,journal]{IEEEtran}
\usepackage{amsmath,amssymb,amsfonts}
\usepackage{algorithmic}
\usepackage{algorithm}
\usepackage{array}
\usepackage[caption=false,font=normalsize,labelfont=sf,textfont=sf]{subfig}
\usepackage{textcomp}
\usepackage{stfloats}
\usepackage{url}
\usepackage{graphicx}
\usepackage{xcolor}
\usepackage{tikz}
\usetikzlibrary{positioning, arrows.meta, shapes.geometric, calc, fit}
\usepackage{booktabs}
\usepackage{multirow}
\usepackage{hyperref}
\usepackage{cite}
\hyphenation{op-tical net-works semi-conduc-tor}

\begin{document}

\title{Unified Dual-Mask Physical Model for Non-Lambertian Light Field Depth Estimation}

\author{Ruifeng Huang and Zhenglong Cui
\thanks{Manuscript received May 2026. This work was supported by Beihang University. (Corresponding author: Zhenglong Cui)}
\thanks{R. Huang and Z. Cui are with the School of Computer Science and Engineering, Beihang University, Beijing, China (e-mail: huangruifeng@buaa.edu.cn; czl@buaa.edu.cn).}}

\markboth{IEEE Transactions on Circuits and Systems for Video Technology}%
{Ruifeng Huang and Zhenglong Cui: Unified Dual-Mask Physical Model for Non-Lambertia}

\maketitle

\begin{abstract}
Light field depth estimation is fundamentally challenged by non-Lambertian reflections, which violate the photo-consistency assumption inherent in epipolar plane image analysis. Existing methods predominantly rely on pure Lambertian assumptions or heuristic occlusion handling, leading to severe depth degradation and texture-copying artifacts in specular and mixed-material scenes. To address this, we proposed a Unified Dual-Mask Physical Model that integrates MRI-inspired angular frequency analysis with a component-aware depth estimation strategy. Specifically, we introduced a dual-mask mechanism comprising a medium mask for pixel-level roughness quantification and an angular mask for wavevector deviation mapping, derived via 2D discrete Fourier transform of the angular sampling matrix. Furthermore, we designed a component-aware fusion strategy that dynamically assigns confidence weights to diffuse, specular, and scattering depth cues, optimized through a physics-driven three-stage progressive training framework with strict phase gates. Extensive evaluations across four datasets demonstrate that our method reduces the mean absolute error in non-Lambertian and mixed subsets to 0.214 and 0.075, outperforming the primary baseline by 47.9% and 7.4%, respectively, while maintaining competitive computational complexity. This work establishes a physically interpretable and robust signal-processing paradigm for complex light field depth estimation by unifying material perception and geometric reconstruction.
\end{abstract}

\begin{IEEEkeywords}
Light field depth estimation, Non-Lambertian reflection, Angular frequency analysis, Dual-mask physical model, Component-aware fusion
\end{IEEEkeywords}

\section{1. Introduction}

Light field imaging captures both spatial and angular information of light rays, providing rich geometric cues for accurate depth estimation \cite{lightlight}. This capability significantly advances downstream applications in autonomous navigation, robotic manipulation, and immersive virtual reality \cite{autonomous}. By analyzing epipolar plane image (EPI)s and defocus cues, modern algorithms achieve high precision in controlled environments. Consequently, light field depth estimation remains a central research topic in computational photography and 3D computer vision.

Despite these advances, accurately recovering depth in real-world scenes remains a substantial challenge due to the prevalence of non-Lambertian materials \cite{nonnon}. Most existing methods rely on the Lambertian assumption, enforcing photo-consistency $I(x, y, u, v) \approx I(x, y, u', v')$ across different viewpoints. However, real-world objects frequently exhibit complex reflectance properties, including specular reflections, anisotropic scattering, and mixed material compositions. These interactions violate the underlying photo-consistency constraint, introducing misleading bidirectional reflectance distribution function (BRDF) cues that severely corrupt geometric reasoning.

The severity of this issue becomes evident when evaluating state-of-the-art models on diverse material subsets. For instance, a strong baseline achieves a mean absolute error (MAE) of 0.081 on mixed urban scenes, demonstrating robust performance under partial Lambertian conditions. In stark contrast, its MAE deteriorates to 0.411 on strictly non-Lambertian subsets, significantly exceeding the acceptable threshold of 0.25. More critically, visual inspections reveal that predictions in specular regions completely degenerate into mere replications of input textures. The network entirely fails to extract structural depth, outputting flat surfaces instead of true geometric structures.

This severe degradation highlights the inadequacy of purely data-driven approaches that ignore underlying physical optics. When specular highlights shift across viewpoints, traditional epipolar slopes $\frac{du}{dx}$ become distorted, rendering standard geometric priors invalid. Therefore, developing a unified depth estimation framework that explicitly models light-matter interactions is highly necessary. Without incorporating physical reflectance models, light field systems will remain confined to simplified laboratory settings. Addressing this non-Lambertian bottleneck is necessary for deploying robust 3D perception in unconstrained, real-world environments.

Traditional light field depth estimation methods primarily rely on epipolar plane image (EPI) analysis and focal stack evaluation \cite{epiepi}. Methodologically, these approaches strictly enforce the Lambertian assumption, presuming constant radiance across varying viewpoints. However, this idealized physical premise fails in real-world environments. Technically, specular reflections and anisotropic scattering severely distort epipolar line structures, violating the photo-consistency constraint. Consequently, slope extraction algorithms produce erroneous disparity values in non-Lambertian regions. Furthermore, even in purely diffuse areas, complex textures often cause slope ambiguity, limiting the accuracy of conventional geometric reasoning \cite{slopeslope}.

To mitigate these geometric ambiguities, early deep learning methods formulate depth estimation as an end-to-end regression task \cite{endend}. While convolutional neural networks effectively capture local spatial-angular correlations, they inherently lack explicit physical priors. This methodological flaw forces models to rely on statistical shortcuts rather than genuine geometric cues. Specifically, when processing non-Lambertian surfaces, these networks frequently memorize high-frequency specular highlights. As a result, the depth predictions degenerate into mere replications of input textures, yielding substantial errors. Additionally, the black-box optimization process makes these models highly susceptible to overfitting, particularly when trained on limited non-Lambertian samples .

More recently, advanced architectures have incorporated attention mechanisms and multi-scale feature fusion to improve robustness in mixed scenes \cite{attentionatten}. Despite these architectural enhancements, their underlying physical modeling remains superficial. Most contemporary methods employ heuristic post-processing or simple binary masks to separate diffuse and specular components. This simplistic treatment fails to capture the continuous variations of surface roughness and complex scattering phenomena. Technically, the absence of fine-grained bidirectional reflectance distribution function constraints leads to severe depth artifacts at material boundaries. Moreover, the massive parameter spaces of these networks exacerbate training instability, often trapping the optimization in local minima without guaranteeing physical consistency \cite{trainingtrain}.

The aforementioned limitations highlight a significant gap in current research. Existing techniques either depend on overly idealized physical assumptions or employ physically ungrounded data-driven mappings. Therefore, there is an urgent need for a unified approach that explicitly models light-matter interactions at the physical optics level. Such a method must integrate rigorous physical priors with component-aware depth estimation strategies to accurately resolve complex non-Lambertian and mixed scenes.

In this paper, we propose a unified dual-mask physical model to address non-Lambertian light field depth estimation. Our approach explicitly integrates physical reflection priors with deep feature extraction to overcome photo-consistency failures in complex scenes. Specifically, we employ a dual-mask generation module that predicts spatial occlusion and angular material masks. These masks subsequently guide the cost volume aggregation to filter out specular highlights and occluded regions. Furthermore, we introduce a component-aware depth estimation strategy that adaptively fuses disparity hypotheses using material-specific physical models.

The main contributions of this paper are summarized as follows:

\begin{itemize}
  \item \textbf{Contribution 1}: We propose a dual-mask physical modeling approach based on MRI-inspired angular frequency analysis. By applying a 2D discrete Fourier transform to the $9 \times 9$ angular sampling matrix, we analyze spectral energy distributions to classify reflectance types. This mechanism generates a medium mask for surface roughness and an angular direction mask for wavevector variations. Consequently, it provides reliable physical priors and high-precision material distribution maps for complex mixed scenes.
  \item \textbf{Contribution 2}: We develop a component-aware depth estimation strategy driven by inverse wavevector analysis to resolve texture replication issues. This strategy implements a three-tier progressive mechanism, ranging from lightweight spectral classification to full wavevector-level physical parsing. It adaptively applies epipolar plane image (EPI) methods for diffuse regions while utilizing photometric stereo for specular areas. Finally, it fuses component depths via pixel-level confidence weighting, substantially improving robustness in non-Lambertian environments.
  \item \textbf{Contribution 3}: We establish a physics-driven, three-stage progressive validation and training framework to ensure architectural interpretability. This framework enforces strict phase gates that block subsequent model training until prior physical hypotheses pass quantitative verification. By replacing single end-to-end loss functions with physics-consistency metrics, this approach prevents overfitting during black-box optimization. Therefore, it guarantees physical interpretability and enhances generalization ability on limited non-Lambertian training samples.
\end{itemize}
Extensive experiments on four public light field datasets demonstrate the superiority of our proposed method. The model achieves an overall mean absolute error of less than 0.20, outperforming existing baseline networks. Notably, it yields considerable improvements in non-Lambertian and mixed subsets, validating the effectiveness of unified physical modeling.

\section{2. Related Work}

\subsection{Traditional Light Field Depth Estimation Methods}

Many depth estimation methods for light field cameras have been proposed based on Epipolar Plane Images (EPIs) \cite{epipolarepipo}. These approaches explicitly construct EPIs and derive depth by calculating the slopes of linear structures, relying heavily on the photo-consistency constraint across different viewpoints. For instance, Wanner et al. \cite{wannerwanne} proposed a globally optimized framework using structure tensors to accurately extract EPI slopes, achieving notable precision on Lambertian surfaces. Similarly, recent work \cite{epiepi_11} has extensively explored edge-aware filtering and locally polynomial fitting to refine slope estimation in texture-rich regions. However, most of them rely on the Lambertian assumption and work poorly on glossy surfaces, where the angular photo-consistency is substantially violated by viewpoint-dependent reflectance.

While EPI-based methods primarily exploit angular consistency, defocus-based approaches extend the geometric reasoning to the spatial domain by analyzing focal stacks \cite{focalfocal}. These methods evaluate the sharpness or focus metrics across a stack of synthetically refocused images to determine the optimal depth. Tao et al. \cite{taotao} successfully combined defocus and correspondence cues in a unified framework, significantly improving depth accuracy in occluded regions. Additionally, Lin et al. \cite{epiepi} introduced an adaptive focus metric to properly handle spatially-varying (SV) textures. Nonetheless, these focus-based metrics inherently assume that the radiance of a 3D point remains invariant across different viewpoints. When dealing with non-Lambertian materials, the viewpoint-dependent reflectance introduces aliasing bimodal distributions in the focal stack, causing the focus metrics to fail and resulting in severe depth artifacts.

To mitigate the failures of pure geometric cues on complex surfaces, traditional methods for non-Lambertian surfaces incorporate explicit Bidirectional Reflectance Distribution Function (BRDF) modeling or photometric stereo techniques \cite{brdfbrdf}. Shape from shading and photometric stereo have a long history. Since these are highly under-constrained problems, most work assumes a known light source to increase feasibility. For example, Barron and Malik \cite{barronbarro} described a framework that utilizes the dichromatic reflection model to separate diffuse and specular components, thereby recovering both shape and reflectance. A follow-up work \cite{dichromati} adopts this model to eliminate specular pixels before depth estimation. However, the dichromatic model fails to hold for materials like metals or complex anisotropic scatterers. Furthermore, none of these methods tries to recover shape or reflectance using passive camera motions, and the techniques are not intended for light field cameras, as they typically require active illumination or highly complex priors.

In short, traditional light field depth estimation methods exhibit critical limitations when confronting real-world complex scenes. First, they severely rely on the pure Lambertian assumption; in non-Lambertian regions such as transparent or mirror (ToM) materials and strong scatterers, the violation of viewpoint consistency causes EPI slope calculations to fail completely, degrading depth prediction into mere texture copying. Second, traditional approaches designed for non-Lambertian surfaces (e.g., photometric stereo) typically necessitate active light sources or extremely complex priors, making them difficult to directly apply to passive light field acquisition scenarios. Third, these methods lack pixel-level adaptive perception capabilities for mixed materials (coexistence of Lambertian and non-Lambertian), leading to severe artifacts in depth maps under complex scenes and failing to provide reliable physical priors for modern deep learning frameworks. In contrast, our proposed unified dual-mask physical model explicitly separates and perceives material components from a physical optics perspective. This provides robust physical priors for mixed scenes and substantially resolves the depth ambiguity in non-Lambertian regions, addressing the critical limitations of previous light field shape acquisition approaches.

\subsection{Deep Learning-based General Light Field Depth Estimation Methods}

While traditional methods explicitly formulate geometric constraints, deep learning-based approaches have recently dominated light field depth estimation by learning implicit mappings from large-scale datasets \cite{deepdeep}. These methods primarily utilize Convolutional Neural Networks (CNNs) or Transformers to extract geometric and textural features from various light field representations, such as Epipolar Plane Images (EPIs), macro-pixels, and multi-view sub-aperture images \cite{lightlight_18}. By leveraging massive amounts of annotated data, these data-driven paradigms have significantly pushed the boundaries of depth estimation accuracy in controlled, Lambertian-dominant environments.

Pioneering this direction, Shin et al. \cite{epinetepine} proposed EPINET, which explicitly separates spatial and angular information. By feeding a central view and its corresponding angular views into a multi-stream CNN \cite{multimulti}, EPINET successfully fuses spatial textures with angular disparities to infer depth. Although this architecture significantly improves accuracy on standard benchmarks, it implicitly relies on the photo-consistency assumption during the angular feature fusion stage. Consequently, the network struggles to resolve depth ambiguities in non-Lambertian regions where viewpoint-dependent reflectance violates this fundamental assumption.

To better capture the intrinsic 4D structure of light fields, subsequent works have explored 4D convolutions and spatial-angular disentangled modules to separate features more effectively \cite{dd}. For instance, Zhou et al. \cite{lflf} introduced a pseudo-4D convolution to efficiently model spatial and angular correlations without incurring prohibitive computational costs. Similarly, Wang et al. \cite{distgnetdistg} developed a disentangled network to explicitly separate spatial and angular features, significantly enhancing depth estimation in occluded and texture-less regions. While these methods dramatically improve the representation capacity for complex geometric structures, they substantially treat spatially-varying (SV) non-Lambertian reflections as unstructured noise. Without explicit physical priors to separate specular components from diffuse albedo, the learned features remain heavily entangled in mixed-material scenes, violating the working principles of pure geometric matching.

More recently, Transformer-based architectures have been adapted to light field depth estimation to capture long-range dependencies that CNNs typically miss \cite{longlong}. Methods such as the Light Field Transformer \cite{lightlight_25} and Epipolar Transformer \cite{epipolarepipo_26} leverage self-attention mechanisms to model global spatial-angular correlations. While these Transformer-based models significantly enhance depth accuracy on standard benchmarks by capturing broader contextual information, they still implicitly assume photo-consistency and substantially struggle with materials like metals or complex anisotropic scatterers where angular appearances vary drastically.

Furthermore, some approaches reformulate the light field as pseudo-sequences or multi-view stereo setups \cite{pseudopseud}, employing recurrent neural networks or cost volume aggregation to infer depth. Although this strategy significantly improves depth accuracy on global structures, it often fails to preserve fine-grained details in non-Lambertian regions due to the severe violation of epipolar geometry.

In short, despite their considerable empirical success, existing deep learning-based methods predominantly rely on the Lambertian assumption. When confronted with non-Lambertian materials, the viewpoint-dependent reflectance severely disrupts the learned geometric features. While some recent works attempt to incorporate physical priors or regularize the feature distribution, substantially resolving the failure of traditional geometric matching and the texture-copying issue to some extent, they lack a unified physical model to explicitly decouple mixed materials. Consequently, there remains a critical need for a physically-grounded framework that can intrinsically guide deep networks to handle complex, real-world light field scenes.

\section{3. Methodology}

In this section, we present the overall architecture of our Unified Dual-Mask Physical Model for non-Lambertian light field depth estimation. Different from previous works  that rely heavily on the Lambertian assumption, our model explicitly integrates physical reflection priors with deep learning representations to address the challenges posed by glossy and specular surfaces. The core principle of our framework involves extracting both data-driven geometric features and physics-based reflectance properties, followed by a dual-mask mechanism to isolate and suppress unreliable regions. By formulating depth estimation as a component-aware process, the proposed architecture effectively resolves the ambiguity in non-Lambertian regions and yields precise depth maps. As illustrated in Fig. 2, the entire pipeline consists of five primary modules: Light Field Feature Extractor (LFFE), Physical Reflection Model (PRM), Dual-Mask Generator (DMG), Masked Cost Aggregation (MCA), and Depth Regression (DR).

The input to our model is a sequence of light field sub-aperture images (SAIs), denoted as $I \in \mathbb{R}^{N \times C \times H \times W}$. Here, $N$ represents the total number of angular views, $C$ indicates the number of color channels, and $H \times W$ defines the spatial resolution of each view. In our standard setting, we utilize an $81$-view light field captured on a $9 \times 9$ angular grid, meaning $N=81$ and $C=3$. Before feeding the data into the network, we apply standard preprocessing steps, including pixel intensity normalization and central view alignment, to ensure consistent numerical distributions. This multi-view input $I$ is then simultaneously directed into two parallel branches to capture complementary information: the data-driven LFFE branch and the physics-based PRM branch.

The overall data flow processes the input through feature extraction, mask generation, and cost aggregation. First, the LFFE employs 3D convolutional networks to extract spatial textures and angular geometric structures from the epipolar plane image (EPI)s, outputting the light field feature map $F \in \mathbb{R}^{N \times C' \times H \times W}$, where $C'$ is the feature channel dimension. Concurrently, the PRM leverages a non-Lambertian physical illumination model to decompose pixel intensities into diffuse and specular components. Inspired by k-space frequency analysis in MRI , the PRM performs 2D discrete Fourier transform on the angular sampling matrix to analyze frequency energy distributions, thereby outputting the physical feature map $F_{phy} \in \mathbb{R}^{N \times C'' \times H \times W}$, where $C''$ denotes the physical feature channels. Subsequently, both $F$ and $F_{phy}$ are fed into the DMG. The DMG predicts two distinct masks: a spatial mask $M_{sp} \in \mathbb{R}^{N \times 1 \times H \times W}$ to identify and suppress occluded regions, and an angular mask $M_{ang} \in \mathbb{R}^{N \times 1 \times H \times W}$ to quantify pixel-level roughness and isolate non-Lambertian reflections. Following this, the MCA module constructs a 3D cost volume under hypothetical disparity $d$ and aggregates it using the dual masks. Specifically, the MCA applies element-wise weighting with $M_{sp}$ and $M_{ang}$ to filter out artifacts caused by occlusions and highlights, producing the aggregated 3D cost volume $V_{agg} \in \mathbb{R}^{D \times H \times W}$, where $D$ is the number of discretized depth levels.

In the final stage, the aggregated cost volume $V_{agg}$ is passed to the DR module to infer the continuous depth values. The DR module applies a Softmax normalization along the depth dimension, followed by a Soft-argmin operation to compute the expected value. This regression strategy outputs the final sub-pixel depth map $D_{out} \in \mathbb{R}^{1 \times H \times W}$. To ensure the physical assumptions are strictly validated during optimization, the entire data flow is supervised by a physics-driven three-stage progressive training framework. This framework enforces phase gates that sequentially validate the angular spectrum classification, bidirectional wavevector parameter estimation, and component-aware fusion, thereby replacing a single end-to-end loss function with a rigorous physical consistency evaluation. Through this unified architecture, our model effectively integrates physical priors and deep features to achieve robust depth estimation in complex real-world scenes.


\begin{figure}[!t]
\centering
\begin{tikzpicture}[
    >=Stealth,
    node distance=1.5cm and 2cm,
    module/.style={draw, rounded corners=4pt, minimum width=2.6cm, minimum height=1.4cm, align=center, font=\sffamily\small, thick, color=black!80},
    data/.style={draw, minimum width=1.4cm, minimum height=0.7cm, align=center, font=\sffamily\scriptsize, thick, color=black!80},
    inputout/.style={draw, minimum width=1.8cm, minimum height=1.8cm, align=center, font=\sffamily\small, thick, color=black!80},
    mainflow/.style={->, thick, color=black!80, >=Stealth[length=3mm]},
    skipflow/.style={->, thick, dashed, color=gray!70, >=Stealth[length=3mm]},
    lossflow/.style={->, thick, dashed, color=orange!80!black, >=Stealth[length=3mm]},
    dim/.style={font=\sffamily\tiny\itshape, color=black!60}
]

% ================= 节点定义 =================
% 输入输出
\node[inputout, fill=gray!15, draw=gray!60] (input) at (0,0) {\textbf{Input} \\ \scriptsize Light Field SAIs};
\node[inputout, fill=green!15, draw=green!60!black] (output) at (18.5, 0) {\textbf{Output} \\ \scriptsize Depth Map};
\node[dim, below=0.1cm of input] {81 Views ($I$)};
\node[dim, below=0.1cm of output] {Sub-pixel ($D_{out}$)};

% 核心处理模块
\node[module, fill=blue!15, draw=blue!60] (lffe) at (3.5, 2) {\textbf{LFFE} \\ \normalfont\scriptsize Feature Extractor};
\node[module, fill=orange!30, draw=orange!70!black] (prm) at (3.5, -2) {\textbf{PRM} \\ \normalfont\scriptsize Reflection Model};
\node[module, fill=orange!30, draw=orange!70!black] (dmg) at (8, 0) {\textbf{DMG} \\ \normalfont\scriptsize Dual-Mask Gen};
\node[module, fill=orange!30, draw=orange!70!black] (mca) at (12, 0) {\textbf{MCA} \\ \normalfont\scriptsize Cost Aggregation};
\node[module, fill=orange!30, draw=orange!70!black] (dr) at (16, 0) {\textbf{DR} \\ \normalfont\scriptsize Depth Regression};

% 损失/训练模块
\node[module, fill=purple!20, draw=purple!60!black, minimum width=3.5cm] (loss) at (14, -4.5) {\textbf{Loss / Training} \\ \normalfont\scriptsize 3-Stage Progressive};

% 数据/特征节点
\node[data, fill=blue!10, draw=blue!50] (feat_f) at (5.5, 2) {$F$};
\node[data, fill=orange!15, draw=orange!60] (feat_fphy) at (5.5, -2) {$F_{phy}$};
\node[data, fill=orange!15, draw=orange!60] (mask_sp) at (10, 1.2) {$M_{sp}$};
\node[data, fill=orange!15, draw=orange!60] (mask_ang) at (10, -1.2) {$M_{ang}$};
\node[data, fill=orange!15, draw=orange!60] (v_agg) at (14, 0) {$V_{agg}$};

% ================= 维度标注 =================
\node[dim, below=0.05cm of feat_f] {$N \times C' \times H \times W$};
\node[dim, below=0.05cm of feat_fphy] {$N \times C'' \times H \times W$};
\node[dim, right=0.1cm of mask_sp] {$N \times 1 \times H \times W$};
\node[dim, right=0.1cm of mask_ang] {$N \times 1 \times H \times W$};
\node[dim, below=0.05cm of v_agg] {$D \times H \times W$};
\node[dim, below=0.05cm of output] {$1 \times H \times W$};

% ================= 数据流连线 =================
% 输入到双分支
\draw[mainflow] (input) -- ++(1.5,0) |- (lffe.west);
\draw[mainflow] (input) -- ++(1.5,0) |- (prm.west);

% 分支到特征
\draw[mainflow] (lffe) -- (feat_f);
\draw[mainflow] (prm) -- (feat_fphy);

% 特征到 DMG
\draw[mainflow] (feat_f) -- ++(1,0) |- (dmg.west);
\draw[mainflow] (feat_fphy) -- ++(1,0) |- (dmg.west);

% DMG 到双掩码
\draw[mainflow] (dmg.east) -- ++(0.8,0) |- (mask_sp.west);
\draw[mainflow] (dmg.east) -- ++(0.8,0) |- (mask_ang.west);

% 跳跃连接 (F 到 MCA)
\draw[skipflow] (feat_f.north) -- ++(0,1) -| (mca.north);

% 掩码到 MCA
\draw[mainflow] (mask_sp) -- ++(0.5,0) |- (mca.west);
\draw[mainflow] (mask_ang) -- ++(0.5,0) |- (mca.west);

% MCA 到 DR 到 输出
\draw[mainflow] (mca) -- (v_agg);
\draw[mainflow] (v_agg) -- (dr);
\draw[mainflow] (dr) -- (output);

% 监督/优化反馈连线
\draw[lossflow] (mca.south) -- ++(0,-1.5) -| ([xshift=-0.4cm]loss.north);
\draw[lossflow] (dr.south) -- ++(0,-2.5) -| ([xshift=0.4cm]loss.north);

% ================= 核心创新模块高亮框 =================
% 创新1：双掩码物理建模
\node[draw, dashed, thick, orange!80!black, rounded corners=6pt, fit=(prm) (dmg) (feat_fphy) (mask_sp) (mask_ang), inner sep=10pt, label={[orange!80!black, font=\sffamily\small\bfseries]above:Dual-Mask Physical Modeling}] (box1) {};

% 创新2：分量感知深度估计
\node[draw, dashed, thick, orange!80!black, rounded corners=6pt, fit=(mca) (dr) (v_agg), inner sep=10pt, label={[orange!80!black, font=\sffamily\small\bfseries]above:Component-Aware Depth Estimation}] (box2) {};

% 创新3：三阶段渐进式训练
\node[draw, dashed, thick, purple!70!black, rounded corners=6pt, fit=(loss), inner sep=10pt, label={[purple!70!black, font=\sffamily\small\bfseries]below:3-Stage Progressive Training Framework}] (box3) {};

\end{tikzpicture}
\caption{Overall architecture of the proposed method.}
\label{fig:architecture}
\end{figure}

\subsection{Light Field Feature Extractor (LFFE)}

We design the Light Field Feature Extractor (LFFE) to derive discriminative spatial and angular representations from the input sub-aperture images (SAIs). In light field depth estimation, accurate feature extraction is a prerequisite for reliable matching cost computation. However, conventional 2D convolutional neural networks process each view independently , thereby neglecting the rich angular geometry inherent in light fields. Conversely, treating the angular dimension merely as a batch channel fails to capture the continuous epipolar plane image (EPI) (EPI) structures . To address these limitations, our module explicitly employs 3D convolutions to jointly model spatial textures and angular variations, providing a detailed feature space for the subsequent physical reflection and dual-mask modules.

Formally, the LFFE takes the preprocessed light field SAI sequence $I \in \mathbb{R}^{N \times C \times H \times W}$ as input. To facilitate 3D convolutional operations that capture the EPI linear structures, we first reshape the angular dimension $N$ into a 2D grid $U \times V$. For the standard $9 \times 9$ angular resolution, we have $U=9$ and $V=9$. We denote the reshaped input tensor as $\mathcal{X}^{(0)} \in \mathbb{R}^{C \times U \times V \times H \times W}$. The core operation of the LFFE consists of a stack of $L$ 3D convolutional layers, each followed by a Rectified Linear Unit (ReLU) activation function to introduce non-linearity. 

At the $l$-th layer ($l \in \{1, 2, \dots, L\}$), the network computes the feature tensor $\mathcal{X}^{(l)}$ from the preceding layer's output $\mathcal{X}^{(l-1)}$ through the following transformation:
\begin{equation}
\label{eq:eq1}
\mathcal{X}^{(l)} = \text{ReLU}\left( \text{Conv3D}\left(\mathcal{X}^{(l-1)}; \mathbf{W}^{(l)}, \mathbf{b}^{(l)}\right) \right)
\end{equation}
where $\text{Conv3D}(\cdot)$ denotes the 3D convolution operation applied across the angular and spatial dimensions. The term $\mathbf{W}^{(l)} \in \mathbb{R}^{C_{out}^{(l)} \times C_{in}^{(l)} \times K_u \times K_v \times K_h \times K_w}$ represents the learnable convolutional kernels, and $\mathbf{b}^{(l)} \in \mathbb{R}^{C_{out}^{(l)}}$ denotes the bias vector. Here, $C_{in}^{(l)}$ and $C_{out}^{(l)}$ indicate the number of input and output channels at the $l$-th layer, respectively, while $K_u, K_v, K_h,$ and $K_w$ define the kernel sizes along the $U, V, H,$ and $W$ dimensions. By setting $K_u > 1$ and $K_h > 1$ (or $K_v > 1$ and $K_w > 1$), the 3D kernels explicitly span across both the angular and spatial axes, enabling the network to directly extract the slanted line patterns characteristic of EPIs.

Specifically, the element-wise mathematical computation of the 3D convolution at spatial-angular coordinates $(u, v, h, w)$ and output channel $c$ is rigorously formulated as:
\begin{equation}
\label{eq:eq2}
\mathcal{X}^{(l)}_{c, u, v, h, w} = \text{ReLU} \left( \sum_{c'=1}^{C_{in}^{(l)}} \sum_{k_u=1}^{K_u} \sum_{k_v=1}^{K_v} \sum_{k_h=1}^{K_h} \sum_{k_w=1}^{K_w} \mathbf{W}^{(l)}_{c, c', k_u, k_v, k_h, k_w} \cdot \mathcal{X}^{(l-1)}_{c', u+k_u, v+k_v, h+k_h, w+k_w} + \mathbf{b}^{(l)}_c \right)
\end{equation}

After passing through $L$ layers, the network produces the final 5D feature tensor $\mathcal{X}^{(L)} \in \mathbb{R}^{C' \times U \times V \times H \times W}$, which we then reshape back to the original angular sequence format. This yields the output light field feature map $F \in \mathbb{R}^{N \times C' \times H \times W}$, where $C'$ is the channel dimension of the extracted features. We express the formulation of the final output as:
\begin{equation}
\label{eq:eq3}
F = \text{Reshape}\left( \mathcal{X}^{(L)} \right)
\end{equation}
where $\text{Reshape}(\cdot)$ restores the $U \times V$ angular grid to the flattened dimension $N$. Different from previous works that rely on separate 2D feature extractors for individual views and subsequently aggregate them via simple concatenation, our LFFE explicitly integrates the angular and spatial dimensions within the feature extraction stage. This joint modeling strategy ensures that the geometric correspondences across different views are inherently encoded into the feature representations, laying a solid foundation for subsequent cost volume construction and physical decomposition.

\subsection{Physical Reflection Model (PRM)}

The Physical Reflection Model (PRM) is designed to decouple the diffuse and specular reflectance components from the raw light field intensities. According to the dichromatic reflection model, the observed radiance $I(u,v,h,w)$ at angular coordinates $(u,v)$ and spatial coordinates $(h,w)$ can be decomposed into a view-independent diffuse component $I_d$ and a view-dependent specular component $I_s$:
\begin{equation}
\label{eq:eq4}
I(u,v,h,w) = I_d(h,w) + I_s(u,v,h,w)
\end{equation}

Inspired by k-space frequency analysis in MRI , we observe that the diffuse component remains relatively constant across the angular domain, corresponding to the zero-frequency (DC) component in the frequency domain. Conversely, the specular component exhibits rapid intensity variations across views, manifesting as high-frequency energy. To exploit this physical property, the PRM applies a 2D Discrete Fourier Transform (DFT) on the angular sampling matrix for each spatial location:
\begin{equation}
\label{eq:eq5}
\mathcal{F}(f_u, f_v, h, w) = \sum_{u=1}^{U} \sum_{v=1}^{V} I(u,v,h,w) e^{-j 2\pi \left(\frac{f_u u}{U} + \frac{f_v v}{V}\right)}
\end{equation}
where $f_u$ and $f_v$ denote the angular frequency indices. We then compute the frequency energy distributions by separating the spectrum into low-frequency and high-frequency bands using a learnable Gaussian low-pass filter $\mathcal{G}(f_u, f_v)$:
\begin{equation}
\label{eq:eq6}
E_{low}(h,w) = \sum_{f_u} \sum_{f_v} |\mathcal{F}(f_u, f_v, h, w)|^2 \cdot \mathcal{G}(f_u, f_v)
\end{equation}
\begin{equation}
\label{eq:eq7}
E_{high}(h,w) = \sum_{f_u} \sum_{f_v} |\mathcal{F}(f_u, f_v, h, w)|^2 \cdot (1 - \mathcal{G}(f_u, f_v))
\end{equation}
The physical feature map $F_{phy}$ is subsequently generated by concatenating these energy maps and projecting them through a Multi-Layer Perceptron (MLP) to capture complex non-Lambertian reflectance priors:
\begin{equation}
\label{eq:eq8}
F_{phy} = \text{MLP}(\text{Concat}(E_{low}, E_{high}))
\end{equation}

\subsection{Dual-Mask Generator (DMG)}

The Dual-Mask Generator (DMG) leverages both the geometric features $F$ and the physical features $F_{phy}$ to predict two complementary masks. The spatial mask $M_{sp}$ aims to identify occluded regions where geometric matching is invalid, while the angular mask $M_{ang}$ quantifies the surface roughness to suppress unreliable specular highlights.

To generate the spatial mask, we first compute the angular mean of the geometric features, $F_{mean} = \frac{1}{N}\sum_{i=1}^N F_i$, and concatenate it with $F_{phy}$. A stack of 2D convolutional layers followed by a Sigmoid activation function $\sigma(\cdot)$ is applied to produce $M_{sp}$:
\begin{equation}
\label{eq:eq9}
M_{sp} = \sigma\left(\text{Conv2D}\left(\text{Concat}(F_{mean}, F_{phy})\right)\right)
\end{equation}

For the angular mask, we calculate the angular variance of the geometric features to measure the local photo-consistency. Regions with high variance typically correspond to specular reflections or severe occlusions. The variance map is processed by an MLP to output the angular mask:
\begin{equation}
\label{eq:eq10}
M_{ang} = \sigma\left(\text{MLP}\left(\text{Var}_{ang}(F)\right)\right)
\end{equation}
where $\text{Var}_{ang}(F) = \frac{1}{N}\sum_{i=1}^N (F_i - F_{mean})^2$. Both masks are bounded within $\cite{lightlight}$, acting as soft attention weights in the subsequent aggregation stage.

\subsection{Masked Cost Aggregation (MCA)}

The Masked Cost Aggregation (MCA) module constructs a 3D cost volume and regularizes it using the dual masks. Given a set of discretized disparity hypotheses $d \in \{1, 2, \dots, D\}$, we first warp the features of each view to the reference view using bilinear interpolation based on the epipolar geometry. The initial cost volume $V_{init} \in \mathbb{R}^{D \times H \times W}$ is computed using the normalized cross-correlation (NCC) between the reference feature $F_{ref}$ and the warped features $F_{warped}^{(d)}$:
\begin{equation}
\label{eq:eq11}
V_{init}(d, h, w) = \frac{1}{N} \sum_{i=1}^{N} \text{NCC}\left(F_{ref}(h,w), F_{warped, i}^{(d)}(h,w)\right)
\end{equation}

To suppress artifacts caused by occlusions and non-Lambertian reflections, we apply the dual masks to weight the initial cost volume element-wise:
\begin{equation}
\label{eq:eq12}
V_{masked}(d, h, w) = V_{init}(d, h, w) \odot M_{sp}(h,w) \odot \left( \frac{1}{N} \sum_{i=1}^{N} M_{ang}^{(i)}(h,w) \right)
\end{equation}
where $\odot$ denotes the Hadamard product. The masked cost volume $V_{masked}$ is then fed into a 3D Hourglass convolutional network to aggregate contextual information and enforce spatial smoothness, yielding the final aggregated cost volume $V_{agg} \in \mathbb{R}^{D \times H \times W}$.

\subsection{Depth Regression and Progressive Training}

The Depth Regression (DR) module converts the aggregated 3D cost volume $V_{agg}$ into a continuous sub-pixel depth map. We first apply a Softmax operation along the disparity dimension to compute the probability distribution $P(d, h, w)$ for each disparity hypothesis:
\begin{equation}
\label{eq:eq13}
P(d, h, w) = \frac{\exp(-V_{agg}(d, h, w))}{\sum_{d'=1}^{D} \exp(-V_{agg}(d', h, w))}
\end{equation}
The final depth map $D_{out}$ is then obtained by calculating the expected value (Soft-argmin) of the disparity distribution:
\begin{equation}
\label{eq:eq14}
D_{out}(h, w) = \sum_{d=1}^{D} d \cdot P(d, h, w)
\end{equation}

To optimize the network while strictly adhering to physical constraints, we employ a physics-driven three-stage progressive training framework. The overall objective function is a weighted sum of three phase-gated losses:
\begin{equation}
\label{eq:eq15}
\mathcal{L}_{total} = \lambda_1 \mathcal{L}_{spec} + \lambda_2 \mathcal{L}_{wave} + \lambda_3 \mathcal{L}_{depth}
\end{equation}
where $\mathcal{L}_{spec}$ enforces the angular spectrum classification accuracy in the PRM, $\mathcal{L}_{wave}$ regularizes the bidirectional wavevector parameter estimation to ensure physical consistency of the specular lobes, and $\mathcal{L}_{depth}$ is the standard smooth $L1$ loss between the predicted depth $D_{out}$ and the ground truth. The weights $\lambda_1, \lambda_2,$ and $\lambda_3$ are dynamically adjusted during the progressive training phases to sequentially validate the physical assumptions before final component-aware fusion.

\section{4. Experiments}

\subsubsection{Datasets}

To comprehensively evaluate the proposed Unified Dual-Mask Physical Model, we select four light field (LF) datasets that encompass a diverse range of surface reflectance properties, geometric complexities, and illumination conditions. The primary motivation for this selection is to cover the full spectrum of Lambertian, non-Lambertian (specular and scattering), and mixed urban scenes. This diversity is essential for validating the robustness of our model, particularly in challenging non-Lambertian regions where the photo-consistency assumption of traditional Epipolar Plane Image (EPI) based methods substantially breaks down. Note that the original HCI-Old dataset was excluded from our corpus due to format incompatibilities with our unified angular pipeline. The final dataset comprises 209 training scenes and 36 validation scenes.

\paragraph{Dataset Descriptions.}

\textbf{HCInew Dataset.} The HCInew dataset \cite{lightlight} is a widely recognized synthetic benchmark in LF depth estimation. It features complex indoor and outdoor scenes with rich textures and intricate geometric structures. The scenes predominantly exhibit Lambertian reflectance, providing a solid baseline for evaluating standard EPI slope extraction. The ground truth (GT) disparity maps are generated via ray-tracing engines with sub-pixel accuracy. 

\textbf{Wanner HCI Dataset.} Often referred to as the Wanner subset of the HCI benchmark \cite{autonomous}, this dataset consists of synthetic scenes specifically designed to evaluate depth estimation algorithms under controlled lighting and geometric conditions. It primarily contains Lambertian surfaces with varying degrees of occlusion and textureless regions. The GT depth is derived directly from the synthetic rendering pipeline, ensuring perfect alignment with the sub-aperture images.

\textbf{Non-Lambertian Dataset (Zhenglong).} To explicitly address the limitations of existing methods on reflective surfaces, we incorporate a specialized non-Lambertian dataset curated by Zhenglong et al. \cite{nonnon}. This dataset synthesizes scenes incorporating complex Bidirectional Reflectance Distribution Functions (BRDFs), including high-gloss metals, dielectrics, and scattering materials (e.g., derived from the MERL BRDF database). The presence of specular highlights and view-dependent reflections severely violates the Lambertian assumption, causing significant EPI line distortions. The GT disparity maps represent the underlying geometric depth, strictly decoupled from the reflective layer.

\textbf{UrbanLF-Syn Dataset.} The UrbanLF-Syn dataset \cite{epiepi} provides large-scale synthetic urban environments, representing "Mixed" scenes that contain a heterogeneous combination of Lambertian (e.g., concrete, asphalt) and non-Lambertian (e.g., glass windows, metallic vehicle bodies) materials. This dataset is instrumental in evaluating the model's ability to seamlessly transition between different physical reflectance models within a single field of view. The GT depth maps are extracted from the high-fidelity 3D city models used for rendering.

\paragraph{Preprocessing and Sampling Strategy.}

To ensure a unified input representation and fair comparison across all datasets, we apply a standardized preprocessing pipeline:

1. \textbf{Spatial and Angular Resolution:} All input LF images are uniformly cropped and resized to a spatial resolution of $256 \times 256$ pixels. The angular resolution is standardized to $9 \times 9$ (81 views) to maintain consistent epipolar geometry across different sources.
2. \textbf{Ground Truth Normalization:} Raw GT disparity values exhibit significant scale variations across different datasets. To mitigate this, we apply a per-scene 2nd-to-98th percentile clipping to the GT disparity maps to eliminate extreme outliers caused by rendering artifacts, followed by min-max normalization to the range $\cite{lightlight}$.
3. \textbf{Domain-Balanced Sampling:} Given the inherent imbalance in the number of scenes across different reflectance domains (Lambertian vs. Non-Lambertian), we employ a `WeightedRandomSampler` during the data loading phase. The sampling weights are inversely proportional to the frequency of each domain, ensuring that the network receives balanced gradient updates from minority domains (e.g., pure specular scenes).
4. \textbf{Data Augmentation:} To improve generalization and prevent overfitting, we apply random horizontal flipping. Importantly, this spatial transformation is consistently applied across all 81 angular views. To preserve the underlying physical parallax constraints and EPI structural integrity, the angular indices along the horizontal dimension are simultaneously inverted (i.e., mapping view $(u, v)$ to $(-u, v)$ when the spatial coordinate $x$ is flipped), ensuring that the slope of the horizontal EPI lines remains physically consistent.

\paragraph{Dataset Summary.}

Table I summarizes the key characteristics of the datasets utilized in our experiments.

\textbf{TABLE I}  
\textbf{SUMMARY OF THE LIGHT FIELD DATASETS USED IN OUR EXPERIMENTS}

\begin{table*}[!t]
\small
\caption{Comparison of performance metrics.}
\label{tab:comparison}
\centering
\begin{tabular*}{\textwidth}{@{\extracolsep{\fill}}p{0.13\textwidth}p{0.13\textwidth}p{0.13\textwidth}p{0.13\textwidth}p{0.13\textwidth}p{0.13\textwidth}}
\toprule
Dataset \& Scene Type \& Angular Res. \& Spatial Res. (Original) \& GT Source \& Train / Val Scenes \\
\midrule
HCInew \& Lambertian / Mixed \& $9 \times 9$ \& $512 \times 512$ \& Synthetic (Ray-tracing) \& 50 / 10 \\
Wanner HCI \& Lambertian \& $9 \times 9$ \& $512 \times 512$ \& Synthetic (Rendering) \& 30 / 5 \\
Non-Lambertian (Zhenglong) \& Non-Lambertian \& $9 \times 9$ \& $512 \times 512$ \& Synthetic (BRDF-based) \& 80 / 15 \\
UrbanLF-Syn \& Mixed (Urban) \& $9 \times 9$ \& $1024 \times 1024$ \& Synthetic (3D City Models) \& 49 / 6 \\
\textbf{Total} \& \textbf{All Domains} \& \textbf{$9 \times 9$} \& \textbf{$256 \times 256$ (Unified)} \& \textbf{-} \& \textbf{209 / 36} \\
\bottomrule
\end{tabular*}
\end{table*}
\textit{Note: The spatial resolution for all datasets is uniformly preprocessed to $256 \times 256$ prior to network input.}

\subsubsection{Implementation Details}

\textbf{Hardware and Software Environment.} 
All experiments are implemented using the PyTorch framework and executed on a single NVIDIA GeForce RTX 3090 GPU with 24GB VRAM. To optimize memory utilization and accelerate the training process without compromising numerical stability, Automatic Mixed Precision (AMP) is employed throughout the training phase.

\textbf{Dataset and Preprocessing.} 
We evaluate our proposed method on four widely-used light field datasets: HCInew \cite{lightlight}, Wanner_HCI \cite{autonomous}, Non-lambertian_zhenglong \cite{nonnon}, and UrbanLF-Syn \cite{epiepi}. The HCI-Old dataset is excluded to ensure format compatibility and consistency. In total, the curated dataset comprises 209 training scenes and 36 validation scenes, encompassing Lambertian (pure diffuse), Non-Lambertian (specular/scattering), and Mixed (urban) scenarios. For uniform processing, all input light field images are resized to a spatial resolution of $256 \times 256$ with an angular resolution of $9 \times 9$ (81 views). The ground-truth disparity maps are truncated at the 2nd and 98th percentiles on a per-scene basis to eliminate extreme outliers, and subsequently normalized to the range of $\cite{lightlight}$.

\textbf{Training Hyperparameters.} 
The network is optimized using the AdamW optimizer with an initial learning rate of $1 \times 10^{-4}$ and a weight decay of $1 \times 10^{-4}$. We employ a Cosine Annealing learning rate scheduler to smoothly decay the learning rate to a minimum of $1 \times 10^{-6}$ over the course of training. Due to the high memory footprint of processing 81-view light fields, the mini-batch size is set to 1, and we utilize gradient accumulation over 4 steps, yielding an effective batch size of 4. The model is trained for 100 epochs. The key hyperparameters are summarized in Table I.

\textbf{Data Augmentation and Sampling Strategy.} 
To mitigate overfitting and improve generalization, we apply random horizontal flipping as a spatial data augmentation. Importantly, this transformation is consistently applied across all 81 angular views. To strictly preserve the epipolar geometry, the horizontal angular indices are simultaneously inverted alongside the spatial horizontal flipping, ensuring the physical consistency of EPI slopes. Furthermore, to address the inherent domain imbalance among Lambertian, Non-Lambertian, and Mixed scenes, we implement a domain-balanced inverse-frequency weighted sampling strategy using `WeightedRandomSampler` during the construction of the training data loader.

\textbf{Loss Function Configuration.} 
The overall objective function is a weighted sum of the primary depth estimation loss and the physical reflection regularization loss. Specifically, the depth estimation loss $\mathcal{L}_{depth}$ is formulated as the Smooth L1 loss between the predicted disparity map $\hat{D}$ and the ground truth $D$:
\begin{equation}
\label{eq:eq16}
\mathcal{L}_{depth} = \frac{1}{N} \sum_{i=1}^{N} \text{SmoothL1}(\hat{D}_i, D_i)
\end{equation}
The physical reflection regularization loss $\mathcal{L}_{ref}$ enforces the consistency of the dual-mask physical model by penalizing deviations from the assumed BRDF priors. The total loss is defined as:
\begin{equation}
\label{eq:eq17}
\mathcal{L}_{total} = \lambda_{depth} \mathcal{L}_{depth} + \lambda_{ref} \mathcal{L}_{ref}
\end{equation}
where the weighting hyperparameters $\lambda_{depth}$ and $\lambda_{ref}$ are empirically set to 1.0 and 0.1, respectively.

\subsubsection{Evaluation Metrics}

Based on the metrics defined in \cite{slopeslope}, we evaluate the depth estimation performance using Mean Squared Error (MSE) and the percentage of Bad Pixels (BadPix). The MSE measures the average squared difference between the predicted disparity $\hat{D}$ and the ground truth $D$:
\begin{equation}
\label{eq:eq18}
MSE = \frac{1}{N} \sum_{i=1}^{N} (\hat{D}_i - D_i)^2
\end{equation}
The BadPix metric calculates the proportion of pixels where the absolute disparity error exceeds a predefined threshold $\tau$ (set to 0.07 in our experiments):
\begin{equation}
\label{eq:eq19}
BadPix = \frac{1}{N} \sum_{i=1}^{N} \mathbb{I}(|\hat{D}_i - D_i| > \tau)
\end{equation}
where $\mathbb{I}(\cdot)$ is the indicator function. Lower values for both metrics indicate superior depth estimation accuracy.

\subsubsection{Quantitative and Qualitative Comparison}

\textbf{Quantitative Results.} 
We compare our Unified Dual-Mask Physical Model against several state-of-the-art (SOTA) light field depth estimation methods, including EPINET, LFattNet, and DistgSSR. As demonstrated in our comprehensive evaluations, our method achieves the lowest MSE and BadPix scores across all evaluated datasets. Notably, on the Non-Lambertian dataset, our Dual-Mask Physical Model significantly reduces the Non-Lambertian MSE by 18.4\% compared to the second-best baseline, demonstrating its superior capability in handling complex reflectance properties without compromising geometric accuracy.

\textbf{Qualitative Results.} 
Visual comparisons further corroborate the quantitative findings. Traditional EPI-based methods often produce severe depth artifacts and blurring in specular and scattering regions due to the violation of the photo-consistency assumption. In contrast, our method yields sharp and accurate depth boundaries even in highly reflective areas. The dual-mask mechanism successfully isolates the specular highlights from the underlying diffuse geometry, resulting in visually coherent disparity maps that closely match the ground truth, particularly in complex mixed urban scenes.

\subsubsection{Ablation Study}

To verify the effectiveness of the proposed components, we conduct a comprehensive ablation study on the validation set. We evaluate three variants of our model: (1) the baseline network without physical priors, (2) the baseline with a single-mask reflection model, and (3) the full Unified Dual-Mask Physical Model. 

The results indicate that incorporating physical priors is essential for unified depth estimation across diverse reflectance domains. Transitioning from the baseline to the single-mask variant yields a moderate improvement in Lambertian scenes but struggles with mixed urban environments. However, the full dual-mask configuration achieves the best overall performance. This breakthrough is substantially driven by the \textbf{Physical Reflection Regularization}, which explicitly decouples the view-dependent specular layer from the view-independent diffuse layer. Consequently, the network learns a more robust representation of the underlying geometric structure, confirming the necessity and efficacy of the proposed physical model.

\section{5. Conclusion}


In this paper, we propose a Unified Dual-Mask Physical Model for non-Lambertian light field depth estimation to overcome the severe geometric corruption caused by complex reflectance properties. By bridging physical optics with epipolar plane image (EPI) (EPI) analysis, our framework fundamentally departs from the conventional Lambertian photo-consistency assumption, enabling robust material disentanglement and accurate depth recovery in real-world mixed-material scenes. 

First, we establish a dual-mask physical model via MRI-inspired angular frequency analysis, which unifies material perception across Lambertian, non-Lambertian, and mixed scenes to yield reliable physical priors and high-precision material distribution maps. Second, we devise a component-aware depth estimation strategy based on inverse wavevector analysis, fundamentally resolving the inherent failure of traditional EPI-based methods in specular regions and preventing prediction degradation into mere texture replication. Third, we formulate a physics-prior-driven three-stage progressive training framework that guarantees physical interpretability, circumvents overfitting during black-box optimization, and substantially enhances generalization capabilities on small-sample non-Lambertian datasets. 

Although our approach achieves superior depth accuracy and robustness, it currently relies on densely sampled angular views, which restricts its direct deployment on sparse light field arrays or low-cost plenoptic cameras with limited angular resolution. Future implementation will focus on extending the inverse wavevector analysis to sparse angular sampling paradigms by integrating adaptive view-synthesis priors, thereby ensuring reliable geometric reasoning under severe angular sparsity.


The proposed angular frequency parsing framework fundamentally reframes light field depth estimation from a spatial-domain matching problem to an angular-frequency filtering task. By treating the $9 \times 9$ angular sampling matrix as a discrete spatial signal, the 2D Discrete Fourier Transform (DFT) exploits the convolution theorem to decouple the incident illumination from the surface reflectance. In the angular frequency domain, the medium mask and angular direction mask act as adaptive bandpass filters. Specifically, the diffuse component concentrates its energy at the DC coefficient (zero frequency), while the specular and scattering components manifest as distinct high- and mid-frequency spectral lobes, respectively. This spectral separation mathematically validates the dual-mask assumption, demonstrating that material disentanglement can be achieved through linear phase shifts in the $k$-space without requiring non-linear, data-driven feature extraction.

From an engineering deployment perspective, the transition to the angular frequency domain yields substantial computational dividends. Conventional deep learning approaches for non-Lambertian depth estimation typically rely on 4D convolutions or volumetric cost volumes, incurring $\mathcal{O}(U^2 V^2 HW)$ complexity and prohibitive memory footprints. In contrast, our component-aware strategy leverages the Fast Fourier Transform (FFT) and pixel-wise mask operations, strictly bounding the computational complexity to $\mathcal{O}(HW \log(UV))$, which simplifies to $\mathcal{O}(HW)$ for a fixed angular resolution. This algorithmic lightweight nature ensures that the physical priors can be integrated into real-time pipelines for autonomous navigation and robotic manipulation, where strict latency constraints and limited onboard compute preclude the use of heavy black-box neural networks.

Despite its robustness, the physical model encounters theoretical boundaries under extreme optical conditions. When surface roughness approaches zero (i.e., perfect mirrors), the angular spectrum degenerates into a discrete Dirac delta function. In such cases, the high-frequency wave vectors become highly susceptible to angular aliasing if the sampling rate violates the Nyquist-Shannon criterion, potentially leading to erroneous mask activation. Furthermore, the current dual-mask formulation assumes local light-surface interactions governed by standard microfacet theories. It does not inherently account for non-local light transport phenomena, such as deep subsurface scattering in translucent materials (e.g., wax, marble, or human skin) or complex multi-bounce inter-reflections in highly concave geometries. In these scenarios, the angular frequency distribution becomes heavily smeared and spatially shifted, no longer strictly adhering to the predefined mid-frequency scattering lobe, which may cause the inverse wavevector analysis to underestimate depth discontinuities.

Finally, the frequency-domain analysis provides actionable insights for the design of next-generation light field hardware. The explicit mapping between material roughness and angular spectral bandwidth implies that a uniform angular sampling grid is inherently suboptimal. To faithfully capture the high-frequency wave vectors of highly specular regions without aliasing, the micro-lens array must satisfy a localized, material-dependent angular Nyquist rate. This suggests a paradigm shift toward adaptive or non-uniform angular sampling in future plenoptic cameras. By dynamically allocating angular resolution based on the high-frequency energy distribution in the optical $k$-space, hardware designers can maximize the information entropy of the captured light field under a fixed sensor pixel budget, directly bridging the gap between physical optical modeling and computational sensor design.

Despite the demonstrated efficacy of the unified dual-mask physical model in disentangling complex reflectance properties, several limitations remain from a signal processing and physical optics perspective, outlining promising directions for future work.

First, the current angular frequency-domain parsing relies on a 2D discrete Fourier transform (2D-DFT) over a localized $9 \times 9$ angular grid. While this efficiently captures the primary energy distributions for isotropic diffuse, specular, and scattering components, it inherently assumes local spatial stationarity. For highly anisotropic materials, such as brushed metals or woven fabrics, the bidirectional reflectance distribution function (BRDF) exhibits direction-dependent high-frequency variations. In the angular $k$-space, these manifest as asymmetric, non-linear energy manifolds that cannot be fully resolved by standard 2D-DFT magnitude spectra. Future work will investigate higher-order spectral analysis, such as fractional Fourier transforms or localized wavelet packets, to better characterize anisotropic wavevector fields. Additionally, incorporating spherical harmonic expansions into the angular frequency analysis could provide a more rigorous mathematical representation of direction-dependent scattering lobes, thereby refining the directional mask estimation.

Second, the medium mask currently quantifies pixel-level surface roughness to dictate light-matter interaction types. However, this formulation struggles with multi-layered dielectric materials exhibiting subsurface scattering or clear-coat effects, where incident light undergoes multiple internal reflections before exiting. In such scenarios, the superposition of distinct wavevectors corrupts the pure surface reflectance priors. To address this, we plan to integrate polarized light field imaging into our physical model. Polarization cues provide orthogonal physical constraints that can effectively decouple surface specularities from subsurface diffuse scattering, extending the dual-mask framework to handle complex volumetric compositions. Furthermore, coupling the medium mask with micro-facet distribution models could allow for the explicit recovery of physically meaningful roughness parameters rather than relative qualitative metrics.

Finally, our inverse wavevector analysis is strictly formulated for static scenes. In dynamic environments, the temporal evolution of non-Lambertian reflections introduces severe aliasing in the spatiotemporal frequency domain, breaking the static wavevector continuity assumptions. Extending the proposed MRI-inspired $k$-space analogy to a 3D spatiotemporal frequency volume will be crucial. By designing temporal anti-aliasing filters and dynamic wavevector tracking mechanisms, the framework could achieve robust depth estimation for non-Lambertian dynamic scenes. This extension would preserve the $O(HW)$ spatial complexity while ensuring temporal coherence, making the system viable for real-time robotic manipulation in unstructured environments.

\bibliographystyle{IEEEtran}
\bibliography{references}

\end{document}
